\documentclass[a4paper,11pt]{article}

\usepackage[utf8]{inputenc}
\usepackage[T1]{fontenc}
%\usepackage[french]{babel}

\usepackage[top=25mm,left=40mm,right=40mm]{geometry}
\usepackage{fourier}
\usepackage[babel]{csquotes}
\usepackage[hyphens]{url}
\usepackage{hyperref}
\hypersetup{
    colorlinks,
    citecolor=black,
    filecolor=black,
    linkcolor=black,
    urlcolor=black
}
\usepackage[normalem]{ulem}

\title{TP2 : Extraction d'information}
\author{Traitement automatique de corpus}
\date{21 octobre 2025}

\begin{document}
\maketitle
\thispagestyle{empty}

Sur la base des éléments méthodologiques et des enseignements techniques présentés lors du cours théorique, il est demandé dans le cadre de ce TP :

\begin{itemize}
\item de se familiariser avec les techniques d'enrichissement (extraction de mots-clés, reconnaissance d'entités nommées, analyse de sentiment...)
\item d'appliquer celles-ci au sous-corpus CAMille fourni sur l'UV
\end{itemize}

Pour ce faire, vous utiliserez différentes librairies Python vues au cours.

\vspace{0.7cm}

Les étapes à mettre en \oe uvre sont les suivantes :

\begin{enumerate}
\item Téléchargez le corpus CAMille (Le Soir) sur l'UV et dézippez les 7936 fichiers texte dans le dossier \texttt {tac/data/txt}
\item Choisissez \underline{une année} au choix entre 1887 et 1970 (par exemple 1899, 1923, 1955, etc.)
\item Modifiez le notebook \texttt{s1\_keywords.ipynb} pour extraire uniquement les mots-clés relatifs à cette année
\item Enrichissez de manière itérative la liste de stopwords et générez un nuage de mots pour l'année choisie avec \texttt{s2\_wordcloud.ipynb} (image à inclure dans le rapport)
\item Modifiez le notebook \texttt{s3\_ner.ipynb} pour trouver les entités nommées principales (personnes, organisations et lieux) pertinentes pour cette année
\item Sélectionnez arbitrairement 10 phrases dans les articles de l'année choisie et analysez-les à l'aide du notebook \texttt{s4\_sentiment.ipynb}, puis résumez les résultats dans un tableau indiquant leur polarité et leur subjectivité
%\item Lisez le notebook \url{s9\_gensim.ipynb} pour comprendre les bases du topic modeling avec Gensim
\item Dans un dossier \enquote{tp2}, créez un nouveau notebook reprenant toutes les étapes précédentes et publiez-le (commit + push) sur GitHub
\item Soumettez sur l'UV un court rapport (2-3 pages, format PDF) présentant brièvement les analyses effectuées et les découvertes relatives à l'année choisie.% et expliquant ce que le topic modeling pourrait apporter de plus
\end{enumerate}

Pour toute question, n'hésitez pas à me contacter via \texttt{max.de.wilde@ulb.be} en mettant systématiquement Denis Lebailly en copie (\texttt{denis.lebailly@ulb.be}).

\end{document}